\subsection{C.3: finite canonical templates in every bounded coefficient degree}
\label{C:canonical-subsection}

Throughout this subsection, $k$ is algebraically closed of characteristic
$p>2$. Fix $m\geq0$, a nonsingular symmetric matrix $J$ of size $m$, a
degree-compatible monomial order, and a total order $\prec$ on the set $k$.
The last order is used only to order eigenvalues; it need not be compatible
with the field operations. Over $k$, this defines the split group $O_m$.
The answer below consists of finite polynomial parameter templates for each
pair $(n,d)$, together with finite operations on their specialized members.
It leaves no matrix-isomorphism quotient in the parameter set.

Here and below a finite algorithm over $k$ uses exact field operations,
equality tests, complete factorization of univariate polynomials over $k$,
and comparisons in the chosen order. Such access is part of the assertion.
The symbolic parameter calculations themselves use polynomial arithmetic,
branching on vanishing of coefficients, and finite elimination. No uniform
Turing presentation of an arbitrary field or of an arbitrary total order is
asserted. A template domain will be a finite Boolean combination of
polynomial equalities and inequalities, together with specified comparisons
between its eigenvalue parameters. Empty domains may be retained; deciding
nonemptiness for an arbitrary additional order is not required.

\begin{lemma}[Proved: effective polynomial projection in the required scope]
\label{C:canonical-polynomial-elimination}
Over an algebraically closed field, every existential formula whose atoms
are polynomial equalities and inequalities can be converted, in finitely
many polynomial operations and coefficient branches, into a quantifier-free
Boolean combination of such atoms. The construction works in every
characteristic. In particular, existence of a matrix satisfying finitely
many polynomial conditions is an explicit polynomial condition on any
remaining parameters.
\end{lemma}
\begin{proof}
It suffices to eliminate one variable $z$ from a conjunction
\[
 P_1(y,z)=\cdots=P_a(y,z)=0,\qquad Q(y,z)\ne0,
\]
where the product of all inequations has been denoted by $Q$; the empty
product is $1$. Disjunctive normal form handles arbitrary Boolean formulas.
Branch on which leading coefficients vanish, and execute the Euclidean
algorithm in $z$ on each branch. Its degrees strictly decrease, and its
divisions use coefficients already declared nonzero. Thus finitely many
branches compute the monic greatest common divisor $G$ of the nonzero
$P_i$, with coefficients rational in $y$ and denominators nonzero on that
branch.

If every $P_i$ is zero after specialization, the desired $z$ exists exactly
when $Q$ is not the zero polynomial: an algebraically closed field is
infinite. This is a finite disjunction of coefficient inequalities. Otherwise
put $D=\max_i\deg_z P_i$, using the initial degree bound. If $\deg G=0$
there is no solution. If $\deg G>0$, a solution exists exactly when
$G$ does not divide $Q^D$. Indeed all roots of $G$ are roots of $Q$ exactly
when each root of $G$, with its multiplicity at most $D$, occurs in $Q^D$.
This argument includes inseparable polynomials. Polynomial division by the
monic $G$ tests divisibility; a nonzero remainder is a finite disjunction
of coefficient inequalities. Clearing the already nonzero denominators
produces polynomial conditions in $y$. Iterate this procedure for all
quantified variables. Complementation also handles negated existential
conditions. Every step is finite, proving the claim.
\end{proof}

\begin{lemma}[Known canonical reduction, with the tuple specialization checked]
\label{C:canonical-tuples}
For fixed positive integers $r,n$, there is a finite, explicitly generated
list of polynomial matrix tuples
\[
 \mathcal T_\alpha(\lambda)
   =(T_{\alpha,1}(\lambda),\ldots,T_{\alpha,r}(\lambda))
       \in M_n(k[\lambda])^r
\]
on polynomial/eigenvalue-order domains $D_\alpha\subseteq k^{s_\alpha}$,
such that their members are exactly one representative of each simultaneous
similarity class of $r$ matrices. A specialized tuple can be reduced to its
unique representative and decomposed into indecomposable simultaneous
blocks by finite exact operations in the access model above.
\end{lemma}
\begin{proof}
Use Belitskii's canonical reduction in the form proved by Sergeichuk:
the reduction is finite and canonical, its parameter domains are computed
from polynomial rank tests and eigenvalue ordering, and fixed matrix size
has finitely many box patterns. His final stabilizer also gives the
canonical direct-sum decomposition
\cite[author version, \S1.4, pp.10--13; Theorems 1.3, 2.2, 2.3,
pp.13, 23--24, 27--28; Remark 2.1, p.28]{C-Sergeichuk2000}.

Here is the explicit input and finite-list construction needed for this
application. Put $N=rn$ and encode a tuple by the $N$-square block matrix
\[
 \mathcal M(T)=
 \begin{pmatrix}
 T_1&T_2&\cdots&T_r\\
 0&0&\cdots&0\\
 \vdots&\vdots&&\vdots\\
 0&0&\cdots&0
 \end{pmatrix},\qquad
 \Lambda=\{\operatorname{diag}(S,\ldots,S):S\in M_n(k)\}.
\]
The algebra $\Lambda$ is a reduced block algebra: its diagonal blocks are
all linked and its off-diagonal blocks are zero. Its invertible elements
act exactly by simultaneous similarity. The subspace of matrices supported
on the first block row is preserved, so reducing an encoded matrix gives
another encoded matrix.

Enumerate the finitely many rectangular box partitions of an $N$-square
array and the finitely many allowed rank and Weyr shapes for each box.
Give every free eigenvalue its own parameter, in order of first occurrence,
and retain the first-row support conditions. Eigenvalues in distinct free
Weyr boxes keep distinct parameter slots even when their specialized values
agree; only copies of one eigenvalue within its Weyr block share a slot.
Include the final stable
scalar cells: each cell $\lambda I$ is a trivial Weyr block whose scalar
has a parameter slot unless an earlier part of the record already
forces its value. A free slot is retained even when its value happens
to be $0$ or $1$. This includes matrices on which the reduction stops
immediately; in particular $n=1$ gives the tuple
$(\lambda_1,\ldots,\lambda_r)$.
Replay the reduction, including this terminal stable-cell record, on each
candidate, retaining it exactly when the reduction fixes the matrix with
the stated record. Linear consistency and rank changes are decided by minors;
choose the first nonzero pivot in a fixed order whenever a choice is needed.
These tests, with denominators cleared on their nonvanishing branches,
give the stated finite domain. This is a construction of the domains,
not just their abstract existence. Every actual input supplies one of the
enumerated patterns, so the list is exhaustive.

Use the deterministic reduction record as the pattern label, omitting
zero-width subdivisions. First-nonzero-pivot conditions make pivot branches
disjoint. Merge repeated descriptions of one record, with the same parameter
slots in their fixed order, by taking the finite union of their domains.
The record is recoverable by running the fixed reduction on its canonical
matrix, including the terminal stable-cell convention just specified.
Every parameter occurs in a specified scalar entry, so its value is
recoverable as well. Canonical uniqueness therefore gives exactly one
indexed representative of each tuple class. In particular no further
orbit equivalence is imposed on the list.

For a specialized tuple, linear elimination and univariate root extraction
execute the reduction. The direct-sum decomposition in Theorem 2.2 is a
simultaneous block decomposition for this repeated-diagonal algebra.
Reading its blocks and reducing them again gives the individual
indecomposable tuples and their multiplicities.
\end{proof}

\begin{theorem}[Proved: explicit canonical parameter templates for rational $O_m$-modules]
\label{C:Om-canonical-templates}
For every $n\geq1$ and $d\geq0$, a finite algorithm constructs a finite list
of polynomial matrix templates and polynomial/eigenvalue-order domains
parametrizing, without repetitions, all indecomposable $n$-dimensional
rational $O_m$-modules of minimal coefficient degree $d$. Their union over
$(n,d)$ is an explicit parameter set $\mathscr I$ for all nonzero
indecomposable rational $O_m$-modules.

Duality is a finite algorithm on $\mathscr I$. The self-dual part is
partitioned by explicit polynomial conditions into symmetric and
alternating types. The remaining parameters occur in computable
two-element duality orbits. There is no unclassified isomorphism quotient
in any of these assertions.
\end{theorem}
\begin{proof}
\emph{The finite coalgebras.}
Write
\[
 H=k[x_{ij}:1\leq i,j\leq m]/(X^{\mathsf t}JX-J).
\]
The determinant relation makes $X$ invertible, and its inverse is
$J^{-1}X^{\mathsf t}J$. Thus this presentation already gives the coordinate
Hopf algebra of $O_m$, with
\[
 \Delta(x_{ij})=\sum_\ell x_{i\ell}\otimes x_{\ell j},\qquad
 \epsilon(x_{ij})=\delta_{ij},\qquad
 \mathsf S(X)=J^{-1}X^{\mathsf t}J .
\]
Compute a Gr\"obner basis in the fixed degree-compatible order. Let $C_d$
be the span of standard monomials of degree at most $d$. Normal reduction
does not increase degree, so these monomials are a basis of the image of
polynomials of degree at most $d$. They give nested finite-dimensional
spaces with union $H$. Comultiplication of a monomial of degree at most
$d$ has degree at most $d$ in each tensor factor. The antipode is linear
on the generators. Consequently
\[
 \Delta(C_d)\subset C_d\otimes C_d,\qquad
 \mathsf S(C_d)=C_d,\qquad C_d C_e\subset C_{d+e}.
                                                        \tag{C.3.T1}
\]
The equality for the antipode uses $\mathsf S^2=1$. Each $C_d$ is therefore
a finite coalgebra with a fixed, explicitly computed, nested basis.

Write that basis as $b_1,\ldots,b_r$ and compute constants
\[
 \Delta(b_c)=\sum_{a,b}\gamma_c^{ab}b_a\otimes b_b,\qquad
 \epsilon(b_c)=e_c,\qquad
 \mathsf S(b_c)=\sum_a s_{ac}b_a .
                                                        \tag{C.3.T2}
\]
The dual algebra $A_d=C_d^*$ has multiplication
$u_a u_b=\sum_c\gamma_c^{ab}u_c$ and unit $\sum_c e_c u_c$.
Its $n$-dimensional modules are precisely the tuples satisfying
\[
 T_aT_b=\sum_c\gamma_c^{ab}T_c\quad(1\leq a,b\leq r),
 \qquad \sum_c e_cT_c=I_n .                              \tag{C.3.T3}
\]
Indeed the associated right coaction is
$v_j\mapsto\sum_{i,a}(T_a)_{ij}v_i\otimes b_a$; its coassociativity and
counit identities are exactly (C.3.T3). Inclusion $C_d\subset H$ identifies
this category fully faithfully with the rational representations whose
matrix coefficients belong to $C_d$. Every finite-dimensional rational
representation belongs to some such category.

\emph{Canonical parameters and indecomposability.}
Apply Lemma~\ref{C:canonical-tuples} with $r=\dim C_d$. On each resulting
template impose the explicit quadratic and unit equations (C.3.T3).
For $d>0$, also require that the matrices attached to monomials of degree
exactly $d$ are not all zero. For $d=0$ impose no extra degree condition.
This is precisely minimal degree: the span of the matrix coefficients is
unchanged by a change of basis. It prevents duplication between distinct
values of $d$.

The centralizer of the tuple is
\[
 E_T=\{Z\in M_n(k):ZT_a=T_aZ\text{ for every }a\}.
                                                        \tag{C.3.T4}
\]
A module is decomposable exactly when $E_T$ contains an idempotent other
than $0$ and $I_n$: projections give one implication, and the image and
kernel of an idempotent give the other. The existence of such an
idempotent is the polynomial existential condition
\[
 \exists Z\quad
 [Z,T_a]=0\ (\forall a),\quad Z^2=Z,\quad Z\ne0,\quad Z\ne I_n.
                                                        \tag{C.3.T5}
\]
Each matrix inequation means that at least one of its finitely many
entries is nonzero. Lemma~\ref{C:canonical-polynomial-elimination} eliminates
the entries of $Z$. Taking the complement gives a finite Boolean
polynomial condition for indecomposability. Equivalently one may use
comprehensive Gr\"obner branches for this polynomial system, retaining
the parameter conditions of each branch: express an inequation $f\ne0$
by $wf-1=0$, and test whether the specialized ideal contains $1$.
Algebraic closedness identifies properness with existence of a solution.
These operations do not
quantify over, or alter, the eigenvalue-order conditions. Intersecting
them with the canonical domains gives the required list. Uniqueness
comes from Lemma~\ref{C:canonical-tuples}, and full faithfulness ensures
that indecomposability here is indecomposability as a rational module.

\emph{Duality and form types.}
By (C.3.T2), the dual tuple is
\[
             T^\vee_a=\sum_c s_{ac}T_c^{\mathsf t}.        \tag{C.3.T6}
\]
This follows by applying the antipode to the transposed coefficient
matrix of the coaction. Canonical reduction computes its index in
$\mathscr I$. Since the antipode preserves the filtration and squares
to the identity, minimal coefficient degree is preserved.

An invariant bilinear form is an intertwiner $B:U\to U^*$, so its matrix
satisfies the linear equations
\[
 BT_a=T^\vee_a B\quad(\forall a).
                                                        \tag{C.3.T7}
\]
Add $B^{\mathsf t}=B$ or $B^{\mathsf t}=-B$. On every finite pivot stratum
in parameter space, Gaussian elimination gives a basis
$B_1,\ldots,B_\ell$ of this solution space, with rational entries and
denominators known to be nonzero. A nondegenerate form of the selected
type exists exactly when
\[
              \det(z_1B_1+\cdots+z_\ell B_\ell)
\]
is not the zero polynomial in $z$. This is a finite disjunction of
nonvanishing coefficient conditions after clearing denominators.
The equivalence uses that $k$ is infinite. Thus both type loci are
explicit polynomial conditions. By the already verified structure
theorem, Proposition~\ref{C:algebraically_closed_WG}, every self-dual
indecomposable has exactly one of these two types, while a non-self-dual
module has neither.

If a form matrix is wanted, its construction is also finite. The
determinant has degree at most $n$ in every variable. Choose $n+1$
distinct elements of $k$ and test the finite Cartesian grid they form.
A nonzero polynomial of these degree bounds cannot vanish on the
entire grid, by induction on the number of variables. The first
nonzero value supplies $B$. Distinct scalars can themselves be obtained
as roots of $z^{p^a}-z$ with $p^a>n$, using the stated root access.
This proves all the asserted constructions.
\end{proof}

\begin{theorem}[Proved: a canonical basis and finite ring and exterior-power algorithms]
\label{C:Om-canonical-WG}
Use the explicit index set $\mathscr I$ of
Theorem~\ref{C:Om-canonical-templates}, and denote its module by $U_i$.
Partition its self-dual indices into $\mathscr I_{\mathrm s}$ and
$\mathscr I_{\mathrm a}$ according to symmetric and alternating type.
Order indices lexicographically by $(d,n,\alpha,\lambda)$, using the fixed
order on $k$ for the parameter vector. Let $\mathscr I_{\mathrm h}$ contain
the smaller index in each non-self-dual pair.

Then $WG_{\mathrm{sp}}(BO_m)$ is the free abelian group on the following
explicitly parametrized basis:
\[
 \begin{array}{c|c|c}
 \text{index}&\text{orthogonal representative}&
       \text{underlying Green-ring class}\\ \hline
 i\in\mathscr I_{\mathrm s}&(U_i,q_i)&[U_i]\\
 i\in\mathscr I_{\mathrm a}&H(U_i)&2[U_i]\\
 i\in\mathscr I_{\mathrm h}&H(U_i)&[U_i]+[U_{i^\vee}]
 \end{array}                                             \tag{C.3.T8}
\]
Here $q_i$ can be constructed by the finite determinant test above.
All product structure constants and every exterior power of every finite
integral linear combination of these basis elements are computable by a
finite algorithm in the stated field-access model. For actual modules of
coefficient degrees $d,e$, tensor products require only $C_{d+e}$;
the $r$th exterior power of a module of degree $d$ requires only $C_{rd}$.
\end{theorem}
\begin{proof}
Proposition~\ref{C:algebraically_closed_WG} applies to these particular
canonical representatives. It proves both that the three lines of
(C.3.T8) form a basis and that forgetting the form is injective. Its
existence and exclusivity statements apply to the rational representation
category: finite-dimensional endomorphism algebras have nilpotent
radical, idempotents split, and finite-dimensional modules have
Krull--Schmidt decompositions. Thus no additional isometry parameter
remains after the canonical module and its type have been specified.
In alternating type, $H(U_i)$ is one generator because its underlying
module is $U_i^{\oplus2}$. These verifications are precisely those in
the previously cited Quebbemann--Scharlau--Schulte theorem, rather than
an assumption of semisimplicity in positive characteristic.

We give the operations explicitly. Represent any basis generator by
its actual underlying coaction matrix $h=(h_{ab})$. For the hyperbolic
generators this is the block sum of a computed tuple and its dual.
If two coaction matrices are $h$ and $g$, their tensor product has entries
$h_{ab}g_{cd}$. Reduce each product in the fixed Gr\"obner basis of $H$,
express it in the basis of $C_{d+e}$, and thereby form its tuple of
$A_{d+e}$-matrices. For the exterior power $\Lambda^r$, the entries of
the coaction matrix are the $r$-rowed minors of $h$. Normal reduction
expresses them in $C_{rd}$. This determinant construction uses no division
by $r!$, including when $p$ divides $r!$. For $r$ larger than the
dimension the module is zero; the zeroth exterior power is the unit.

Apply Lemma~\ref{C:canonical-tuples} to the resulting tuple and read its
finite Krull--Schmidt decomposition. For every indecomposable block,
find the least $d'$ for which its coefficient matrices outside $C_{d'}$
vanish, restrict to that nested basis, and canonicalize again. This step
is necessary because the ambient bound need not be the minimal degree
of a summand. It yields its unique index in $\mathscr I$.

Suppose the resulting Green-ring class is
$\sum_{i\in\mathscr I}a_i[U_i]$, with only finitely many nonzero $a_i$.
It came from an actual invariant nondegenerate symmetric form, since
tensor and exterior powers preserve such forms. Consequently
$a_i$ is even for $i\in\mathscr I_{\mathrm a}$ and
$a_i=a_{i^\vee}$ for every non-self-dual pair, by
Proposition~\ref{C:algebraically_closed_WG}. The coefficients in
(C.3.T8) are therefore exactly
\[
 a_i\quad(i\in\mathscr I_{\mathrm s}),\qquad
 a_i/2\quad(i\in\mathscr I_{\mathrm a}),\qquad
 a_i\quad(i\in\mathscr I_{\mathrm h}).                     \tag{C.3.T9}
\]
Injectivity of forgetting the form proves that these integers are the
required structure constants, not merely invariants of the answer.

Products of virtual elements follow by bilinearity. For exterior powers
use the orthogonal-sum identity
$\lambda_t(x+y)=\lambda_t(x)\lambda_t(y)$ and the formal inverse
$\lambda_t(-x)=\lambda_t(x)^{-1}$. Every coefficient of fixed degree
requires only finitely many of the computed actual exterior powers and
ring products. These are the given exterior operations; no special
$\lambda$-ring identity in positive characteristic has been used.
All the indicated steps terminate, proving the theorem.
\end{proof}

\begin{proposition}[Proved: an exact $O_3$ illustration in characteristics $3$ and $5$]
\label{C:Om-canonical-example}
Let $V=k^3$ be the standard $O_3$-module with Gram matrix $I_3$, and
write $\delta=\det(V)$. Both $V$ and $\delta\otimes V$ are absolutely
indecomposable of symmetric type, and
\[
 \Lambda^2 V\simeq\delta\otimes V,\qquad
 \Lambda^3 V\simeq\delta .
                                                        \tag{C.3.T10}
\]
In characteristic $3$, $S=\operatorname{Sym}^2 V$ is an absolutely
indecomposable symmetric-type module of dimension $6$. In characteristic
$5$ it decomposes as $\mathbf1\oplus W$, where $W$ is absolutely
indecomposable of symmetric type and dimension $5$. In the basis of
Theorem~\ref{C:Om-canonical-WG}, using the corresponding named indices,
\[
 [V]^2=
 \begin{cases}
 [S]+[\delta\otimes V],&p=3,\\
 1+[W]+[\delta\otimes V],&p=5,
 \end{cases}
 \qquad
 \lambda^2[V]=[\delta\otimes V],\quad
 \lambda^3[V]=[\delta].                                 \tag{C.3.T11}
\]
\end{proposition}
\begin{proof}
Diagonal sign changes force an endomorphism of $V$ to be diagonal,
and permutation matrices force all diagonal entries to agree.
Hence $\operatorname{End}_{O_3}(V)=k$. The Hodge map identifies the
exterior square with $\delta\otimes V$, and the determinant identifies
the top exterior power with $\delta$. These identities hold as
rational representations of the group scheme.

Identify $S$ with symmetric $3$-square matrices, with action
$A\mapsto gAg^{\mathsf t}$. Put $D_i=E_{ii}$ and
$F_{ij}=E_{ij}+E_{ji}$. Commutation with diagonal sign changes and
permutations forces every equivariant endomorphism $T$ to have the form
\[
 T(D_i)=aD_i+b\sum_{j\ne i}D_j,\qquad T(F_{ij})=cF_{ij}.
\]
The three off-diagonal sign characters are distinct, which justifies
the latter assertion in every odd characteristic. Take
\[
 g=\begin{pmatrix}u&v&0\\-v&u&0\\0&0&1\end{pmatrix},
 \qquad u^2+v^2=1,\quad uv\ne0 .
\]
Such $u,v$ exist over the infinite algebraically closed field $k$.
Comparing the $F_{12}$ coefficients in $T(gD_1g^{\mathsf t})$ and
$gT(D_1)g^{\mathsf t}$ gives $c=a-b$. Conversely the identity and
$K(A)=\operatorname{tr}(A)I_3$ are equivariant. Thus
\[
 \operatorname{End}_{O_3}(S)=kI\oplus kK,\qquad K^2=3K.    \tag{C.3.T12}
\]
The two displayed generators are linearly independent.

For $p=3$ this endomorphism algebra is $k[z]/(z^2)$, which has no
nontrivial idempotent, so $S$ is indecomposable. For $p=5$, $K/3$
is the projection onto $kI_3$, and its complementary summand is
$W=\ker(\operatorname{tr})$. The algebra (C.3.T12) is $k\times k$;
each corner algebra is $k$, so the two summands are indecomposable.
These assertions are absolute: endomorphisms are kernels of the
finite coefficient-commutator system, which commutes with extension
of the coefficient field, and the displayed algebras remain local
on each stated indecomposable summand after any such extension.

The invariant form $\operatorname{tr}(AB)$ on symmetric matrices
has diagonal Gram entries $1,1,1,2,2,2$ in the indicated basis.
It is nondegenerate in both characteristics. For $p=5$ the line
$kI_3$ is nondegenerate, of norm $3$, and its orthogonal complement
is exactly $W$. Since $2$ is invertible, the tensor square decomposes
as its symmetric and alternating eigenspaces. The alternating
eigenspace has the exterior-square form up to the nonzero scalar
$2$; over $k$ that scalar is a square. Alternatively injectivity of
forgetting forms in Proposition~\ref{C:algebraically_closed_WG} gives
the same identities. This proves (C.3.T11).

There is also an executable certificate for the finite calculations:
\path{code/C_finish_coalgebra_canonical.py}, with output
\path{logs/C_finish_coalgebra_canonical.json}. It computes the
coordinate quotient $\mathbb F_p[x_{ij}]/(X^{\mathsf t}X-I)$ itself,
rather than only the group of finite-field points. For each of
$p=3,5$ its degree-compatible Gr\"obner basis has $27$ elements;
$\dim C_d$ for $d=0,1,2,3$ is $1,10,44,119$. The $C_2$
coproduct has $400$ nonzero terms. Per prime, the certificate checks
$44$ coassociativity identities, $88$ counit identities, $88$
antipode identities, $18$ orthogonality identities, $7744$ dual-algebra
module products, $4$ module-unit identities, $176$ invariant-form
coefficient equations, and $10$ Hodge/determinant equations. Its
full centralizer computations give dimensions $1,2,1$ for
$V,S,\Lambda^2 V$ and verify (C.3.T12). These are bounded checks of
the stated example and coalgebra construction; the finite symbolic
canonical-template algorithm is proved above, not claimed to be
implemented by this certificate.
\end{proof}

\paragraph{Scope of this determination.}
For each $(n,d)$, the recipe produces finitely many explicitly testable
parameter templates, and each requested product or exterior power uses
only finitely many of them. It does not claim a finite classification
at fixed $n$ without a coefficient-degree bound, a finite generating
set for the ring, or efficient complexity. It also does not require
the group to have finite representation type. Wildness is compatible
with this parameter-template answer. For fields which are not
algebraically closed, residue division algebras and their hermitian
forms still occur in Proposition~\ref{C:field_WG_structure}; they have
not been classified by this construction. The unrestricted
nonfield-base part of C.3 remains separate.
